%--------------------------------------------------------------------
\section{几何关系 API}
\label{sec:geometric_relations_api}

\subsection{一般原则}

在 \tkzNamePack{tkz-elements} 中，几何对象提供一个统一的方法
\[
\code{object:position(other\_object[, EPS])}
\]
用于判定几何关系。

判定结果取决于第二个参数（\code{other\_object}）的类型。

\medskip
所有位置判定均带有容差。
若未指定，可选参数 \code{EPS} 默认为全局值
\tkzVar{tkz}{tkz.epsilon}。

\medskip
\textbf{约定：}
所有判定结果均以\emph{大写符号字符串}返回，
供 Lua 代码中做逻辑比较使用。

%--------------------------------------------------------------------
\subsection{点与对象}

当被检测的对象是一个点时，结果描述该点相对于某个区域（或边界）的归属关系。

\medskip
可能的返回值：

\begin{center}
\begin{tabular}{ll}
\code{"IN"}  & 严格位于区域内部， \\
\code{"ON"}  & 位于边界上（在容差范围内）， \\
\code{"OUT"} & 严格位于区域外部。
\end{tabular}
\end{center}

\medskip
示例：
\begin{itemize}
\item \code{circle:position(point)} \quad（\code{"IN"}、\code{"ON"}、\code{"OUT"}）
\item \code{triangle:position(point)} \quad（\code{"IN"}、\code{"ON"}、\code{"OUT"}）
\item \code{conic:position(point)} \quad（\code{"IN"}、\code{"ON"}、\code{"OUT"}）
\item \code{line:position(point)} \quad（\code{"ON"} 或 \code{"OUT"}）
\end{itemize}

%--------------------------------------------------------------------
\subsection{对象与对象}

当两个参数都是几何对象时，结果描述它们之间的相互几何关系。

\subsubsection*{直线与直线}

\begin{center}
\begin{tabular}{ll}
\code{"INTERSECT"} & 两直线交于一点， \\
\code{"PARALLEL"}  & 不同的平行直线， \\
\code{"IDENTICAL"} & 重合的直线。
\end{tabular}
\end{center}

\subsubsection*{直线与圆}

\begin{center}
\begin{tabular}{ll}
\code{"DISJOINT"} & 无交点， \\
\code{"TANGENT"}  & 恰有一个公共点， \\
\code{"SECANT"}   & 两个交点。
\end{tabular}
\end{center}

\subsubsection*{圆与圆}

\begin{center}
\begin{tabular}{ll}
\code{"DISJOINT\_EXT"} & 两圆外离， \\
\code{"TANGENT\_EXT"}  & 两圆外切， \\
\code{"SECANT"}        & 两个交点， \\
\code{"TANGENT\_INT"}  & 两圆内切， \\
\code{"DISJOINT\_INT"} & 一个圆严格位于另一个圆内部， \\
\code{"CONCENTRIC"}    & 圆心相同、半径不同， \\
\code{"IDENTICAL"}     & 两圆重合。
\end{tabular}
\end{center}

%--------------------------------------------------------------------
\subsection{三角形与圆锥曲线}

对于 \tkzClass{triangle} 和 \tkzClass{conic}，方法 \code{position()}
目前仅支持传入点作为参数。

\subsubsection*{三角形与点}

\begin{center}
\begin{tabular}{ll}
\code{"IN"}  & 严格位于三角形内部， \\
\code{"ON"}  & 位于边上或顶点上， \\
\code{"OUT"} & 位于三角形外部。
\end{tabular}
\end{center}

\subsubsection*{圆锥曲线与点}

\begin{center}
\begin{tabular}{ll}
\code{"IN"}  & 位于圆锥曲线所围区域内部（按其类型而定）， \\
\code{"ON"}  & 位于圆锥曲线上， \\
\code{"OUT"} & 位于相应区域外部。
\end{tabular}
\end{center}

\medskip
\texttt{注意：}对于双曲线和抛物线，\code{"IN"} 与
\code{"OUT"} 的含义取决于该圆锥曲线类型所关联的区域及其
构造数据（焦点/准线或等价表示）。

%--------------------------------------------------------------------
\subsection{设计理念}

\begin{itemize}
\item 用单一的多态方法集中处理几何关系判定。
\item 大写符号结果保证了清晰性与稳定性。
\item 数值容差提高了边界情形附近的稳健性。
\item 该 API 可扩展：今后可能支持更多对象类型。
\item 若早期版本中曾存在旧的布尔方法，后续可能提供向后兼容的封装。
\end{itemize}
